\section{How to Prepare Mathematical Notes}

A mathematical note is not merely a place where an answer is recorded. Its
purpose is to show how a problem is understood, how relevant information is
selected, how a method is chosen, and how a conclusion follows from
definitions, assumptions, and earlier results.

Mathematical expressions are compact. A single formula may contain several
ideas at once: objects, operations, conditions, and logical relationships.
Reading a formula therefore means more than pronouncing its symbols. A student
should be able to explain what the objects represent, what operation is being
performed, why the expression is meaningful, and what information the formula
provides.

The same principle applies to writing. A sequence of formulas without
explanatory sentences rarely communicates a complete mathematical argument.
Formulas should be connected by words that explain their role. A good solution
tells the reader why a calculation is being performed, which definition or
theorem is being used, whether its assumptions are satisfied, and how the
result answers the original question.

\subsection{Make the Note Self-Contained}

Begin by writing the problem or a complete and faithful formulation of it. The
reader should not need to search through another document to discover what is
being solved. After several weeks, you should still be able to open the note
and understand both the question and the solution.

Next, identify what is given and what must be found, proved, or explained.
Introduce the relevant notation and determine what kinds of mathematical
objects appear in the problem. Before calculating, ask what information is
available and what would count as a satisfactory answer.

\subsection{Make the Reasoning Visible}

Present the principal idea or method before, or while, carrying out the
calculation. Important transitions should be explained with complete
sentences. For example:
\begin{itemize}
    \item Because the determinant is non-zero, the matrix is invertible.
    \item We may therefore multiply both sides by the inverse matrix.
    \item By the definition of the derivative, we consider the following limit.
    \item Since the direction vector is non-zero, it determines a line.
    \item Substituting the result into the original equation verifies the
    solution.
\end{itemize}

Such sentences are not decoration around the mathematics. They are part of the
mathematics because they show the logical relationship between successive
formulas.

It is not necessary to describe every elementary arithmetic operation. The
amount of explanation should be proportional to the problem. A simple exercise
may require only a few sentences, while a more substantial problem may require
a longer argument. The goal is not to make every note long. The goal is to make
every essential step understandable and justified.

\begin{quote}
Write as briefly as possible, but as completely as necessary.
\end{quote}

\subsection{State and Examine the Conclusion}

Finish with a clear answer to the original question. When appropriate, verify
the result by substitution, check that all assumptions and domain restrictions
are satisfied, or interpret the answer geometrically or conceptually.

A result without its context is not a complete solution. The reader should
know what has been established and why it follows from the work presented
above.

\subsection{Using AI to Prepare Notes}

AI may help explain notation, identify missing steps, suggest relevant
definitions, check calculations, or improve the organisation of a solution. It
may also help transform rough work into clear mathematical prose.

However, copying a short AI-generated answer, an unexplained sequence of
formulas, or a final result does not demonstrate mathematical understanding.
Such work will be treated as an incomplete solution and should be revised.

A polished AI-generated note is also insufficient if the student cannot read
its formulas, explain its terminology, justify its transitions, or present the
argument independently. The student is responsible for every statement and
formula included in the submitted work.

When using AI, ask it not only to solve the problem, but also to help you
understand the solution:
\begin{itemize}
    \item What does this formula mean?
    \item What mathematical object does each symbol represent?
    \item Why is this step valid?
    \item Which definition or theorem is being used?
    \item Are the assumptions of the theorem satisfied?
    \item Is there a simpler way to explain this transition?
    \item How can the result be checked?
    \item What should I be able to explain during class?
\end{itemize}

The final note should be written for a human reader. A complete solution should
make the following five stages visible:

\begin{center}
\begin{tikzpicture}[
    stage/.style={
        draw=SectionBlue,
        rounded corners=3pt,
        fill=AccentBlue!35,
        text width=10.2cm,
        minimum height=1.05cm,
        align=center,
        inner sep=6pt
    },
    centralstage/.style={
        stage,
        draw=ExampleGreen,
        fill=ExampleGreenBack
    },
    flowarrow/.style={
        -{Stealth[length=2.6mm]},
        very thick,
        draw=SubsectionBlue
    }
]
\node[stage] (understand) at (0,0)
    {\textbf{1. Understand the problem}\\
     Record the problem, identify the given information, and state the goal.};

\node[stage] (method) at (0,-1.65)
    {\textbf{2. Choose a method}\\
     Select the definitions, theorems, or techniques that connect the data
     with the goal.};

\node[centralstage] (argument) at (0,-3.30)
    {\textbf{3. Construct the argument}\\
     Carry out the necessary steps and connect formulas with sentences that
     explain why each important transition is valid.};

\node[stage] (verify) at (0,-4.95)
    {\textbf{4. Verify the result}\\
     Check calculations, assumptions, domain restrictions, or geometric
     meaning whenever appropriate.};

\node[stage] (conclude) at (0,-6.60)
    {\textbf{5. State the conclusion}\\
     Answer the original question clearly and in its proper context.};

\draw[flowarrow] (understand) -- (method);
\draw[flowarrow] (method) -- (argument);
\draw[flowarrow] (argument) -- (verify);
\draw[flowarrow] (verify) -- (conclude);
\end{tikzpicture}
\end{center}

Preparing such notes is part of learning mathematics. It trains the ability to
read notation, recognise mathematical objects, organise information, construct
arguments, and communicate ideas precisely. These skills are more important
than producing a large number of unexplained answers.
